\IfFileExists{stacks-project.cls}{%
\documentclass{stacks-project}
}{%
\documentclass{amsart}
}

% For dealing with references we use the comment environment
\usepackage{verbatim}
\newenvironment{reference}{\comment}{\endcomment}
%\newenvironment{reference}{}{}
\newenvironment{slogan}{\comment}{\endcomment}
\newenvironment{history}{\comment}{\endcomment}

% For commutative diagrams we use Xy-pic
\usepackage[all]{xy}

% We use 2cell for 2-commutative diagrams.
\xyoption{2cell}
\UseAllTwocells

% We use multicol for the list of chapters between chapters
\usepackage{multicol}

% This is generally recommended for better output
\usepackage{lmodern}
\usepackage[T1]{fontenc}

% For cross-file-references
\usepackage{xr-hyper}

% Package for hypertext links:
\usepackage{hyperref}

% For any local file, say "hello.tex" you want to link to please
% use \externaldocument[hello-]{hello}
\externaldocument[introduction-]{introduction}
\externaldocument[conventions-]{conventions}
\externaldocument[sets-]{sets}
\externaldocument[categories-]{categories}
\externaldocument[topology-]{topology}
\externaldocument[sheaves-]{sheaves}
\externaldocument[sites-]{sites}
\externaldocument[stacks-]{stacks}
\externaldocument[fields-]{fields}
\externaldocument[algebra-]{algebra}
\externaldocument[brauer-]{brauer}
\externaldocument[homology-]{homology}
\externaldocument[derived-]{derived}
\externaldocument[simplicial-]{simplicial}
\externaldocument[more-algebra-]{more-algebra}
\externaldocument[smoothing-]{smoothing}
\externaldocument[modules-]{modules}
\externaldocument[sites-modules-]{sites-modules}
\externaldocument[injectives-]{injectives}
\externaldocument[cohomology-]{cohomology}
\externaldocument[sites-cohomology-]{sites-cohomology}
\externaldocument[dga-]{dga}
\externaldocument[dpa-]{dpa}
\externaldocument[sdga-]{sdga}
\externaldocument[hypercovering-]{hypercovering}
\externaldocument[schemes-]{schemes}
\externaldocument[constructions-]{constructions}
\externaldocument[properties-]{properties}
\externaldocument[morphisms-]{morphisms}
\externaldocument[coherent-]{coherent}
\externaldocument[divisors-]{divisors}
\externaldocument[limits-]{limits}
\externaldocument[varieties-]{varieties}
\externaldocument[topologies-]{topologies}
\externaldocument[descent-]{descent}
\externaldocument[perfect-]{perfect}
\externaldocument[more-morphisms-]{more-morphisms}
\externaldocument[flat-]{flat}
\externaldocument[groupoids-]{groupoids}
\externaldocument[more-groupoids-]{more-groupoids}
\externaldocument[etale-]{etale}
\externaldocument[chow-]{chow}
\externaldocument[intersection-]{intersection}
\externaldocument[pic-]{pic}
\externaldocument[weil-]{weil}
\externaldocument[adequate-]{adequate}
\externaldocument[dualizing-]{dualizing}
\externaldocument[duality-]{duality}
\externaldocument[discriminant-]{discriminant}
\externaldocument[derham-]{derham}
\externaldocument[local-cohomology-]{local-cohomology}
\externaldocument[algebraization-]{algebraization}
\externaldocument[curves-]{curves}
\externaldocument[resolve-]{resolve}
\externaldocument[models-]{models}
\externaldocument[functors-]{functors}
\externaldocument[equiv-]{equiv}
\externaldocument[pione-]{pione}
\externaldocument[etale-cohomology-]{etale-cohomology}
\externaldocument[proetale-]{proetale}
\externaldocument[relative-cycles-]{relative-cycles}
\externaldocument[more-etale-]{more-etale}
\externaldocument[trace-]{trace}
\externaldocument[crystalline-]{crystalline}
\externaldocument[spaces-]{spaces}
\externaldocument[spaces-properties-]{spaces-properties}
\externaldocument[spaces-morphisms-]{spaces-morphisms}
\externaldocument[decent-spaces-]{decent-spaces}
\externaldocument[spaces-cohomology-]{spaces-cohomology}
\externaldocument[spaces-limits-]{spaces-limits}
\externaldocument[spaces-divisors-]{spaces-divisors}
\externaldocument[spaces-over-fields-]{spaces-over-fields}
\externaldocument[spaces-topologies-]{spaces-topologies}
\externaldocument[spaces-descent-]{spaces-descent}
\externaldocument[spaces-perfect-]{spaces-perfect}
\externaldocument[spaces-more-morphisms-]{spaces-more-morphisms}
\externaldocument[spaces-flat-]{spaces-flat}
\externaldocument[spaces-groupoids-]{spaces-groupoids}
\externaldocument[spaces-more-groupoids-]{spaces-more-groupoids}
\externaldocument[bootstrap-]{bootstrap}
\externaldocument[spaces-pushouts-]{spaces-pushouts}
\externaldocument[spaces-chow-]{spaces-chow}
\externaldocument[groupoids-quotients-]{groupoids-quotients}
\externaldocument[spaces-more-cohomology-]{spaces-more-cohomology}
\externaldocument[spaces-simplicial-]{spaces-simplicial}
\externaldocument[spaces-duality-]{spaces-duality}
\externaldocument[formal-spaces-]{formal-spaces}
\externaldocument[restricted-]{restricted}
\externaldocument[spaces-resolve-]{spaces-resolve}
\externaldocument[formal-defos-]{formal-defos}
\externaldocument[defos-]{defos}
\externaldocument[cotangent-]{cotangent}
\externaldocument[examples-defos-]{examples-defos}
\externaldocument[algebraic-]{algebraic}
\externaldocument[examples-stacks-]{examples-stacks}
\externaldocument[stacks-sheaves-]{stacks-sheaves}
\externaldocument[criteria-]{criteria}
\externaldocument[artin-]{artin}
\externaldocument[quot-]{quot}
\externaldocument[stacks-properties-]{stacks-properties}
\externaldocument[stacks-morphisms-]{stacks-morphisms}
\externaldocument[stacks-limits-]{stacks-limits}
\externaldocument[stacks-cohomology-]{stacks-cohomology}
\externaldocument[stacks-perfect-]{stacks-perfect}
\externaldocument[stacks-introduction-]{stacks-introduction}
\externaldocument[stacks-more-morphisms-]{stacks-more-morphisms}
\externaldocument[stacks-geometry-]{stacks-geometry}
\externaldocument[moduli-]{moduli}
\externaldocument[moduli-curves-]{moduli-curves}
\externaldocument[examples-]{examples}
\externaldocument[exercises-]{exercises}
\externaldocument[guide-]{guide}
\externaldocument[desirables-]{desirables}
\externaldocument[coding-]{coding}
\externaldocument[obsolete-]{obsolete}
\externaldocument[fdl-]{fdl}
\externaldocument[index-]{index}

% Theorem environments.
%
\theoremstyle{plain}
\newtheorem{theorem}[subsection]{Theorem}
\newtheorem{proposition}[subsection]{Proposition}
\newtheorem{lemma}[subsection]{Lemma}

\theoremstyle{definition}
\newtheorem{definition}[subsection]{Definition}
\newtheorem{example}[subsection]{Example}
\newtheorem{exercise}[subsection]{Exercise}
\newtheorem{situation}[subsection]{Situation}

\theoremstyle{remark}
\newtheorem{remark}[subsection]{Remark}
\newtheorem{remarks}[subsection]{Remarks}

\numberwithin{equation}{subsection}

% Macros
%
\def\lim{\mathop{\mathrm{lim}}\nolimits}
\def\colim{\mathop{\mathrm{colim}}\nolimits}
\def\Spec{\mathop{\mathrm{Spec}}}
\def\Hom{\mathop{\mathrm{Hom}}\nolimits}
\def\Ext{\mathop{\mathrm{Ext}}\nolimits}
\def\SheafHom{\mathop{\mathcal{H}\!\mathit{om}}\nolimits}
\def\SheafExt{\mathop{\mathcal{E}\!\mathit{xt}}\nolimits}
\def\Sch{\mathit{Sch}}
\def\Mor{\mathop{\mathrm{Mor}}\nolimits}
\def\Ob{\mathop{\mathrm{Ob}}\nolimits}
\def\Sh{\mathop{\mathit{Sh}}\nolimits}
\def\NL{\mathop{N\!L}\nolimits}
\def\CH{\mathop{\mathrm{CH}}\nolimits}
\def\proetale{{pro\text{-}\acute{e}tale}}
\def\etale{{\acute{e}tale}}
\def\QCoh{\mathit{QCoh}}
\def\Ker{\mathop{\mathrm{Ker}}}
\def\Im{\mathop{\mathrm{Im}}}
\def\Coker{\mathop{\mathrm{Coker}}}
\def\Coim{\mathop{\mathrm{Coim}}}

% Boxtimes
%
\DeclareMathSymbol{\boxtimes}{\mathbin}{AMSa}{"02}

%
% Macros for moduli stacks/spaces
%
\def\QCohstack{\mathcal{QC}\!\mathit{oh}}
\def\Cohstack{\mathcal{C}\!\mathit{oh}}
\def\Spacesstack{\mathcal{S}\!\mathit{paces}}
\def\Quotfunctor{\mathrm{Quot}}
\def\Hilbfunctor{\mathrm{Hilb}}
\def\Curvesstack{\mathcal{C}\!\mathit{urves}}
\def\Polarizedstack{\mathcal{P}\!\mathit{olarized}}
\def\Complexesstack{\mathcal{C}\!\mathit{omplexes}}
% \Pic is the operator that assigns to X its picard group, usage \Pic(X)
% \Picardstack_{X/B} denotes the Picard stack of X over B
% \Picardfunctor_{X/B} denotes the Picard functor of X over B
\def\Pic{\mathop{\mathrm{Pic}}\nolimits}
\def\Picardstack{\mathcal{P}\!\mathit{ic}}
\def\Picardfunctor{\mathrm{Pic}}
\def\Deformationcategory{\mathcal{D}\!\mathit{ef}}
\usepackage{cleveref}
\crefname{originallemma}{original lemma}{original lemmas}
\Crefname{originallemma}{Original lemma}{Original lemmas}
\makeatletter
\newlabel{lemma-construction-base-change}{{0H6K}{source}{}{sourcepointer-1}{}}
\newlabel{lemma-construction-base-change@cref}{{[originallemma][0][]0H6K}{[1][0][]source}{}{}{}}
\makeatother
\crefname{originallemma}{original lemma}{original lemmas}
\Crefname{originallemma}{Original lemma}{Original lemmas}
\makeatletter
\newlabel{lemma-construction-opens}{{0H6J}{source}{}{sourcepointer-2}{}}
\newlabel{lemma-construction-opens@cref}{{[originallemma][0][]0H6J}{[1][0][]source}{}{}{}}
\makeatother
\crefname{originallemma}{original lemma}{original lemmas}
\Crefname{originallemma}{Original lemma}{Original lemmas}
\makeatletter
\newlabel{lemma-specialization-flat-pullback}{{0H4L}{source}{}{sourcepointer-3}{}}
\newlabel{lemma-specialization-flat-pullback@cref}{{[originallemma][0][]0H4L}{[1][0][]source}{}{}{}}
\makeatother
\crefname{originallemma}{original lemma}{original lemmas}
\Crefname{originallemma}{Original lemma}{Original lemmas}
\makeatletter
\newlabel{lemma-get-cycles}{{0H59}{source}{}{sourcepointer-4}{}}
\newlabel{lemma-get-cycles@cref}{{[originallemma][0][]0H59}{[1][0][]source}{}{}{}}
\makeatother
\crefname{originalequation}{original equation}{original equations}
\Crefname{originalequation}{Original equation}{Original equations}
\makeatletter
\newlabel{equation-cycle-classes}{{0H58}{source}{}{sourcepointer-5}{}}
\newlabel{equation-cycle-classes@cref}{{[originalequation][0][]0H58}{[1][0][]source}{}{}{}}
\makeatother
\crefname{originallemma}{original lemma}{original lemmas}
\Crefname{originallemma}{Original lemma}{Original lemmas}
\makeatletter
\newlabel{chow-lemma-envelope}{{0GU8}{source}{}{sourcepointer-6}{}}
\newlabel{chow-lemma-envelope@cref}{{[originallemma][0][]0GU8}{[1][0][]source}{}{}{}}
\makeatother
\crefname{originallemma}{original lemma}{original lemmas}
\Crefname{originallemma}{Original lemma}{Original lemmas}
\makeatletter
\newlabel{lemma-relative-cycle-functoriality}{{0H51}{source}{}{sourcepointer-7}{}}
\newlabel{lemma-relative-cycle-functoriality@cref}{{[originallemma][0][]0H51}{[1][0][]source}{}{}{}}
\makeatother
\crefname{originallemma}{original lemma}{original lemmas}
\Crefname{originallemma}{Original lemma}{Original lemmas}
\makeatletter
\newlabel{lemma-uniqueness-extension}{{0H55}{source}{}{sourcepointer-8}{}}
\newlabel{lemma-uniqueness-extension@cref}{{[originallemma][0][]0H55}{[1][0][]source}{}{}{}}
\makeatother
\crefname{originallemma}{original lemma}{original lemmas}
\Crefname{originallemma}{Original lemma}{Original lemmas}
\makeatletter
\newlabel{lemma-construction-bilinear}{{0H6I}{source}{}{sourcepointer-9}{}}
\newlabel{lemma-construction-bilinear@cref}{{[originallemma][0][]0H6I}{[1][0][]source}{}{}{}}
\makeatother
\crefname{originallemma}{original lemma}{original lemmas}
\Crefname{originallemma}{Original lemma}{Original lemmas}
\makeatletter
\newlabel{lemma-construction-push-pull}{{0H6N}{source}{}{sourcepointer-10}{}}
\newlabel{lemma-construction-push-pull@cref}{{[originallemma][0][]0H6N}{[1][0][]source}{}{}{}}
\makeatother
\crefname{originallemma}{original lemma}{original lemmas}
\Crefname{originallemma}{Original lemma}{Original lemmas}
\makeatletter
\newlabel{lemma-get-cycles-dvr}{{0H5A}{source}{}{sourcepointer-11}{}}
\newlabel{lemma-get-cycles-dvr@cref}{{[originallemma][0][]0H5A}{[1][0][]source}{}{}{}}
\makeatother
\crefname{originallemma}{original lemma}{original lemmas}
\Crefname{originallemma}{Original lemma}{Original lemmas}
\makeatletter
\newlabel{lemma-construction-closed}{{0H6M}{source}{}{sourcepointer-12}{}}
\newlabel{lemma-construction-closed@cref}{{[originallemma][0][]0H6M}{[1][0][]source}{}{}{}}
\makeatother
\crefname{originallemma}{original lemma}{original lemmas}
\Crefname{originallemma}{Original lemma}{Original lemmas}
\makeatletter
\newlabel{lemma-family-associated-closed-specialization}{{0H57}{source}{}{sourcepointer-13}{}}
\newlabel{lemma-family-associated-closed-specialization@cref}{{[originallemma][0][]0H57}{[1][0][]source}{}{}{}}
\makeatother
\title{Composition of relative cycles is a relative cycle}
\author{The Stacks Project Authors}
\date{stacks-project:relative-cycles.tex:local-snapshot-20261012; GNU FDL 1.2 or later}
\begin{document}
\maketitle
\setcounter{subsection}{1}
\noindent\textit{Editorial note.} All delivered mathematical text is the author's own, copied from the pinned Stacks Project chapter file relative-cycles.tex; no statement or proof step was rewritten, reordered or removed. The counted span is the author environment and its own proof environment. References inside the span are handed to the assembler as navigation stubs that carry the canonical tag of the cited tag only; whether each of them is load-bearing is recorded on the stub itself.\par\medskip
\noindent\textit{Allowed standard prerequisites (editorial).}\begin{itemize}
\item Base-change compatibility of the composition of families of cycles (relative-cycles chapter, Lemma lemma-construction-base-change, canonical tag 0H6K). cited at proof lines 3009-3010. Reduces the specialization check to the case S the spectrum of a discrete valuation ring.
\item Locality of the composition on opens and compatibility of specialization maps with flat pullback (relative-cycles chapter, Lemma lemma-construction-opens, canonical tag 0H6J, and Lemma lemma-specialization-flat-pullback, canonical tag 0H4L). cited at proof lines 3014-3018. The constructions commute with restriction, so X and Y may be assumed affine.
\item Proper completely decomposed cover putting a relative cycle into the image of the cycle-classes map (relative-cycles chapter, Lemma lemma-get-cycles, canonical tag 0H59) and the cycle-classes construction (relative-cycles chapter, Equation equation-cycle-classes, canonical tag 0H58). cited at proof lines 3020-3023. Provides the completely decomposed proper morphism g : Y' to Y with g\textasciicircum{}*alpha in the image.
\item Envelope decomposition of a prime cycle along a completely decomposed morphism (Chow Homology, Lemma chow-lemma-envelope, canonical tag 0GU8). cited at proof lines 3026-3029. Writes beta\_eta as a sum of pushforwards g\_\{eta,*\}beta'\_eta along the fibres of the cover.
\item Functoriality of relative cycles under pushforward and the uniqueness of extensions agreeing on a dense open (relative-cycles chapter, Lemma lemma-relative-cycle-functoriality, canonical tag 0H51, and Lemma lemma-uniqueness-extension, canonical tag 0H55). cited at proof lines 3032-3035 and 3044-3045. The pieces beta' assemble to a relative cycle whose pushforward equals beta, and relative cycles are detected by their generic fibres; functoriality also reduces the target to the pullback composition.
\item Bilinearity of the composition and its push-pull formula (relative-cycles chapter, Lemma lemma-construction-bilinear, canonical tag 0H6I, and Lemma lemma-construction-push-pull, canonical tag 0H6N). cited at proof lines 3036-3037 and 3042-3043. Linearity reduces to a single summand; push-pull moves the pushforward f\_* outside the composition.
\item Discrete valuation ring version of the cover lemma (relative-cycles chapter, Lemma lemma-get-cycles-dvr, canonical tag 0H5A). cited at proof lines 3046-3047. Further reduces the pulled-back composition to a single cycle class of a flat closed subscheme.
\item Closed-pair formula for the composition and the specialization lemma for families associated to closed subschemes (relative-cycles chapter, Lemma lemma-construction-closed, canonical tag 0H6M, and Lemma lemma-family-associated-closed-specialization, canonical tag 0H57). cited at proof lines 3054-3060. Computes alpha composition beta as the cycle class of the flat fibred product and certifies it is a relative cycle.
\end{itemize}\medskip
\ifdefined\hypertarget\hypertarget{sourcepointer-1}{}\fi
\noindent\textit{Source reference.} Original lemma 0H6K: sources/stacks-chow/tags.txt, source lines 21425--21425. Cited at proof lines 3009-3010: reduces to the discrete valuation ring case.\par\medskip
\ifdefined\hypertarget\hypertarget{sourcepointer-2}{}\fi
\noindent\textit{Source reference.} Original lemma 0H6J: sources/stacks-chow/tags.txt, source lines 21424--21424. Cited at proof lines 3015-3016: lets the proof work locally on Y and X.\par\medskip
\ifdefined\hypertarget\hypertarget{sourcepointer-3}{}\fi
\noindent\textit{Source reference.} Original lemma 0H4L: sources/stacks-chow/tags.txt, source lines 21356--21356. Cited at proof lines 3016-3018: the specialization construction commutes with restriction.\par\medskip
\ifdefined\hypertarget\hypertarget{sourcepointer-4}{}\fi
\noindent\textit{Source reference.} Original lemma 0H59: sources/stacks-chow/tags.txt, source lines 21379--21379. Cited at proof lines 3020-3023: provides the cover g with g\textasciicircum{}*alpha in the image.\par\medskip
\ifdefined\hypertarget\hypertarget{sourcepointer-5}{}\fi
\noindent\textit{Source reference.} Original equation 0H58: sources/stacks-chow/tags.txt, source lines 21378--21378. Cited at proof lines 3022-3023: the image condition on g\textasciicircum{}*alpha.\par\medskip
\ifdefined\hypertarget\hypertarget{sourcepointer-6}{}\fi
\noindent\textit{Source reference.} Original lemma 0GU8: sources/stacks-chow/tags.txt, source lines 20993--20993. Cited at proof lines 3026-3029: decomposes beta\_eta into pushforwards along the fibres of the cover.\par\medskip
\ifdefined\hypertarget\hypertarget{sourcepointer-7}{}\fi
\noindent\textit{Source reference.} Original lemma 0H51: sources/stacks-chow/tags.txt, source lines 21371--21371. Cited at proof lines 3032-3033 and 3044-3045: g\_*beta' and beta are relative cycles, and the target reduces to the pullback composition.\par\medskip
\ifdefined\hypertarget\hypertarget{sourcepointer-8}{}\fi
\noindent\textit{Source reference.} Original lemma 0H55: sources/stacks-chow/tags.txt, source lines 21375--21375. Cited at proof line 3035: g\_*beta' equals beta because they agree at eta.\par\medskip
\ifdefined\hypertarget\hypertarget{sourcepointer-9}{}\fi
\noindent\textit{Source reference.} Original lemma 0H6I: sources/stacks-chow/tags.txt, source lines 21423--21423. Cited at proof lines 3036-3037: reduces to a single summand of the envelope decomposition.\par\medskip
\ifdefined\hypertarget\hypertarget{sourcepointer-10}{}\fi
\noindent\textit{Source reference.} Original lemma 0H6N: sources/stacks-chow/tags.txt, source lines 21428--21428. Cited at proof lines 3042-3043: moves the pushforward f\_* outside the composition.\par\medskip
\ifdefined\hypertarget\hypertarget{sourcepointer-11}{}\fi
\noindent\textit{Source reference.} Original lemma 0H5A: sources/stacks-chow/tags.txt, source lines 21380--21380. Cited at proof lines 3046-3047: reduces to a single flat cycle class.\par\medskip
\ifdefined\hypertarget\hypertarget{sourcepointer-12}{}\fi
\noindent\textit{Source reference.} Original lemma 0H6M: sources/stacks-chow/tags.txt, source lines 21427--21427. Cited at proof lines 3054-3057: computes the composition explicitly.\par\medskip
\ifdefined\hypertarget\hypertarget{sourcepointer-13}{}\fi
\noindent\textit{Source reference.} Original lemma 0H57: sources/stacks-chow/tags.txt, source lines 21377--21377. Cited at proof lines 3059-3060: certifies the computed class is a relative cycle.\par\medskip
\begin{lemma}
\label{lemma-well-defined}
If $\alpha$ and $\beta$ are relative cycles, then so is $\alpha \circ \beta$.
\end{lemma}

\noindent
Let $S$ be a locally Noetherian scheme. Let $X \to Y$
be a morphism of schemes locally of finite type over $S$.
We are going to define a map
$$
z(X/Y, r) \otimes_\mathbf{Z} z(Y/S, e) \longrightarrow z(X/S, r + e),\quad
\alpha \otimes \beta \longmapsto \alpha \circ \beta
$$

\begin{proof}
The formation of $\alpha \circ \beta$ is compatible with base change
by Lemma \ref{lemma-construction-base-change}. Thus we may assume
$S$ is the spectrum of a discrete valuation ring with generic point
$\eta$ and closed point $0$ and we have to show that
$sp_{X/S}((\alpha \circ \beta)_\eta) = (\alpha \circ \beta)_0$.
Since we are trying to prove an equality of cycles, we may
work locally on $Y$ and $X$ (this uses
Lemmas \ref{lemma-construction-opens} and
\ref{lemma-specialization-flat-pullback}
to see that the constructions commute with restriction).
Thus we may assume $X$ and $Y$ are affine.
By Lemma \ref{lemma-get-cycles}
we can find a completely decomposed proper morphism
$g : Y' \to Y$ such that $g^*\alpha$ is in the image of
(\ref{equation-cycle-classes}).

\medskip\noindent
Since the family of morphisms $g_\eta : Y'_\eta \to Y_\eta$
is completely decomposed, we can find $\beta'_\eta \in Z_e(Y'_\eta)$ such that
$\beta_\eta = \sum g_{\eta, *}\beta'_\eta$, see
Chow Homology, Lemma \ref{chow-lemma-envelope}.
Set $\beta'_0 = sp_{Y'/S}(\beta'_\eta)$ so that
$\beta' = (\beta'_\eta, \beta'_0)$ is a relative $e$-cycle on $Y'/S$. Then
$g_*\beta'$ and $\beta$ are relative $e$-cycles
on $Y/S$ (Lemma \ref{lemma-relative-cycle-functoriality})
which have the same value at $\eta$ and hence are equal
(Lemma \ref{lemma-uniqueness-extension}). By linearity
(Lemma \ref{lemma-construction-bilinear})
it suffices to show that $\alpha \circ g_*\beta'$
is a relative $(r + e)$-cycle.

\medskip\noindent
Set $X' = X \times_Y Y'$ and denote $f : X' \to X$ the projection.
By Lemma \ref{lemma-construction-push-pull} we see that
$\alpha \circ g_*\beta' = f_*(g^*\alpha \circ \beta')$.
By Lemma \ref{lemma-relative-cycle-functoriality}
it suffices to show that $g^*\alpha \circ \beta'$
is a relative $(r + e)$-cycle. Using Lemma \ref{lemma-get-cycles-dvr}
and bilinearity this reduces us to the case discussed in the next paragraph.

\medskip\noindent
Assume $\alpha = [Z/X/Y]_r$ and $\beta = [W/Y/S]$
where $Z \subset X$ is a closed subscheme flat and of relative
dimension $\leq r$ over $Y$ and $W \subset Y$ is a closed
subscheme flat and of relative dimension $\leq e$ over $S$.
By Lemma \ref{lemma-construction-closed} we see that
$$
\alpha \circ \beta = [Z \times_X W/X/S]_{r + e}
$$
and $Z \times_X W \subset X$ is a closed subscheme flat over $S$
of relative dimension $\leq r + e$. This is a relative $(r + e)$-cycle by
Lemma \ref{lemma-family-associated-closed-specialization}.
\end{proof}

\end{document}
